\renewcommand{\bibname}{References and source map}
\begin{thebibliography}{99}
\addcontentsline{toc}{chapter}{References and source map}
\raggedright
% Bibliography entries for chapters 01--07. Input inside the shared
% thebibliography environment; these keys are deliberately namespaced.
\bibitem{intro-boe} McLeay, Michael, Amar Radia and Ryland Thomas (14 March 2014).
\textit{Money creation in the modern economy}. Bank of England Quarterly
Bulletin, 54(1), 14--27. Institutional description of modern UK money creation,
including lending constraints; no current quantitative estimate used.
\url{https://www.bankofengland.co.uk/quarterly-bulletin/2014/q1/money-creation-in-the-modern-economy}.

\bibitem{intro-ets} European Commission, Directorate-General for Climate Action.
\textit{About the EU ETS}. Maintained institutional overview, ``How does the
EU ETS work?'' Accessed 19 September 2026. Used only for the cap-and-trade,
monitoring, reporting and allowance-surrender mechanism; no price, effectiveness
estimate or proposed reform treated as an adopted rule.
\url{https://climate.ec.europa.eu/areas-action/carbon-markets/about-eu-ets_en}.

\bibitem{intro-water} Food and Agriculture Organization of the United Nations.
\textit{Water scarcity}. Land, Soil and Water overview; physical, access and
infrastructure scarcity. Accessed 19 September 2026. No numerical scarcity
estimate used.
\url{https://www.fao.org/land-water/water/water-scarcity/en}.

\bibitem{intro-health} World Health Organization (6 May 2025).
\textit{World report on social determinants of health equity}.
ISBN 978-92-4-010758-8. Official publication overview used for the qualitative
relationship of social conditions and resource access to health inequities;
no effect size or individual medical recommendation inferred.
\url{https://www.who.int/publications/i/item/9789240107588}.

\bibitem{intro-ipcc} IPCC (2023). \textit{Climate Change 2023: Synthesis Report.
Summary for Policymakers}. In H. Lee and J. Romero (eds.). Section A.1.1 and
note 5: observed 2011--2020 global warming relative to 1850--1900. Section
B.1.1 and notes 28--31: conditional 2081--2100 warming under named SSP
scenarios. These are dated assessment values, not current-year measurements
or assigned probabilities of future social pathways.
\url{https://www.ipcc.ch/report/ar6/syr/summary-for-policymakers/}.

\bibitem{intro-ipbes} IPBES (2019). \textit{Summary for policymakers of the
global assessment report on biodiversity and ecosystem services}.
S. D\'iaz et al. (eds.). Section A5: the assessment estimate of roughly one
million threatened animal and plant species. See also IPBES Secretariat
(30 May 2019), \textit{A million threatened species? Thirteen questions and
answers}, which explains the estimate. Threatened is not equivalent to
already extinct, and the estimate is not a current census.
\url{https://www.ipbes.net/zh/node/35313}.

\bibitem{intro-resources} UNEP and International Resource Panel
(1 March 2024). \textit{Global Resources Outlook 2024}. Official report
overview: the conditional increase in global resource extraction by about
60 percent between 2020 and 2060 without urgent concerted action.
\url{https://www.unep.org/resources/Global-Resource-Outlook-2024}.

\bibitem{intro-homeostasis} McFarland, Jenny, Mary Pat Wenderoth,
Joel Michael, William Cliff, Ann Wright and Harold Modell (2016).
\textit{A conceptual framework for homeostasis: development and validation}.
Advances in Physiology Education, 40(2), 213--222.
\url{https://doi.org/10.1152/advan.00103.2015}.
This educational framework addresses organism-level homeostasis through
negative feedback; it is not a theory of economic or ecosystem control.

\bibitem{onsager-one} Lars Onsager. \emph{Reciprocal Relations in Irreversible Processes. I.} Physical Review 37, 405--426 (1931). \href{https://doi.org/10.1103/PhysRev.37.405}{doi:10.1103/PhysRev.37.405}. Primary source for the historical reciprocal-relations context, not a validation of EBU.
\bibitem{onsager-two} Lars Onsager. \emph{Reciprocal Relations in Irreversible Processes. II.} Physical Review 38, 2265--2279 (1931). \href{https://doi.org/10.1103/PhysRev.38.2265}{doi:10.1103/PhysRev.38.2265}. Linear force--rate relationships, microscopic reversibility and dissipation-function context; EBU's controller remains a separately declared model.
\bibitem{gaussian} Energy Balance Project. \emph{Local Gaussian EBU: programme reconciliation}. Repository document \texttt{LOCAL\_GAUSSIAN\_EBU\_PROGRAMME\_RECONCILIATION.md}, starting revision \texttt{05f3cb0}, September 2026. Sections 3--8 define the prospective domain, potential, finite identities, common-path interpretation and capacity proposal. This is a prospective model and conditional derivation, not an executed study.
\bibitem{foundation} Energy Balance Project. \emph{V3.0 Local EBU Foundation Draft}. Repository source \texttt{V3.0\_LOCAL\_EBU\_FOUNDATION\_DRAFT.md}. Sections 6--8 supply scoped finite, local, path and sequential identities; Section 10 distinguishes simultaneous group value from candidate allocations. Historical model and execution restrictions remain attached to their statements.
\bibitem{bridge} Energy Balance Project. \emph{The Sequential--Parallel Bridge of EBU}, analytical checkpoint v0.2. Repository source \texttt{SEQUENTIAL\_PARALLEL\_BRIDGE.md}. Sections 3--7 distinguish group transitions, sequential telescoping, endpoint equivalence and comparator-relative interaction. Prospective operational proposals in the document are not scientific execution records.
\bibitem{original} Konrad Grzyb. \emph{The Energy Balance Project, Part I: Physical Intuition and the EBU Idea}. Unified explanatory edition, August 2026, 296 pages. Author-supplied PDF; the content and visual baseline for this full-length revised edition. Original Parts II and III are preserved with their original chapter numbering and historical model scope.
\bibitem{gate1dc} Energy Balance Project. \emph{V3.0 Gate 1D-C: Finalized Execution Manifest}. Repository record \texttt{results/v3.0/gate1dc/MANIFEST.md}, finalized historical evidence. Sections 1--4 record study status, execution identity, source locks and the inventory of thirty registered runs. This manuscript cites that limited historical fact and does not transfer its findings to the Gaussian programme.
\end{thebibliography}

\section*{How to follow a mathematical claim}
The social and environmental references support the bounded claims attached to
them in the opening chapters, not a claim that EBU succeeds. The project
sources identify definitions, conditional mathematics and historical records.
Source versions, PDF hashes and the full revision disposition accompany the
editable manuscript.

The local-to-global argument in this edition is the cancellation of unchanged
factors. Its Gaussian examples use a newly declared quadratic; they do not
reinterpret historical threshold fixtures. The exact finite expansion is
derived directly from that quadratic. The path sum uses linearity, the chain
rule and the fundamental theorem of calculus. The sequence and closed-cycle
identities use consecutive endpoint cancellation.

The field definition, the path convention and any proposed capacity rule have
different roles. A proof about the sum of path contributions cannot supply a
missing ownership map. Likewise, a physical conservation equation cannot by
itself choose which consequences belong in a potential. These boundaries are
part of the claim, including when a source file is committed in Git.
